\documentclass[12pt]{article}

\usepackage{amsmath}
\usepackage{amssymb}
\usepackage{amsthm}
\usepackage[headheight=15pt]{geometry}
\usepackage{booktabs}
\usepackage{array}
\usepackage{hyperref}
\usepackage{microtype}
\usepackage{parskip}
\usepackage{titlesec}
\usepackage{enumitem}
\usepackage{fancyhdr}
\usepackage{xcolor}
\usepackage{listings}

\geometry{
  letterpaper,
  top=1.25in, bottom=1.25in, headheight=15pt,
  left=1.25in, right=1.25in
}

\pagestyle{fancy}
\fancyhf{}
\rhead{\small PhiBase Arithmetic}
\lhead{\small Draft}
\cfoot{\thepage}

% Theorem environments
\newtheorem{theorem}{Theorem}[section]
\newtheorem{lemma}[theorem]{Lemma}
\newtheorem{corollary}[theorem]{Corollary}
\theoremstyle{definition}
\newtheorem{definition}[theorem]{Definition}
\theoremstyle{remark}
\newtheorem{remark}[theorem]{Remark}

% Shorthand
% \varphi is already defined
\newcommand{\Z}{\mathbb{Z}}
\newcommand{\R}{\mathbb{R}}
\newcommand{\Zphi}{\mathbb{Z}[\varphi]}
\newcommand{\Phib}[2]{P(#1,\,#2)}
\newcommand{\norm}[1]{\mathcal{N}(#1)}

% Code style
\lstset{
  basicstyle=\small\ttfamily,
  columns=flexible,
  keepspaces=true
}

% ─────────────────────────────────────────────────────────────────────────────
\begin{document}

\title{%
  \textbf{PhiBase Arithmetic and the Digital Slide Rule}\\[6pt]
  \large Exact Geometric Computation in $\mathbb{Z}[\varphi]$,\\
  Fibonacci-Indexed Lookup Tables,\\
  and a $\varphi$-Native GPU Crystal Lattice
}

\author{%
  SpaceStruts Research\\[4pt]
  \small\texttt{spacestruts.com}
}

\date{\today}

\maketitle
\thispagestyle{fancy}

% ─────────────────────────────────────────────────────────────────────────────
\begin{abstract}
We describe a lookup-based arithmetic system grounded in the mathematics
of the golden ratio $\varphi = (1+\sqrt{5})/2$.  Every scalar in the
system is written in \emph{PhiBase notation} as $P(p,n) = p\varphi + n$,
where $p$ and $n$ are integers.  Two precomputed tables support all
arithmetic operations without floating-point arithmetic at runtime.
The \emph{Fibonacci Spine} (Table~2) encodes exact powers of $\varphi$
indexed by a single integer $k$; multiplication reduces to index addition.
The \emph{Dense Grid} (Table~1) covers $\approx 13{,}000$ values in the
positive cone at $10^{-8}$ precision; comparison reduces to integer rank
comparison.  Together the tables constitute a \emph{digital slide rule}
operating on a Fibonacci scale rather than a logarithmic one.

We then apply this foundation to dodecahedral geometry.  The six
basis vectors of the icosahedral lattice are expressed exactly in
$\Zphi$, and we prove an \emph{integrality theorem} for reflections:
reflection through any of the six dodecahedral face-normal planes maps
$\Zphi^{3}$ into $\tfrac{1}{5}\Zphi^{3}$, with denominator \emph{exactly
$5$} in all cases.  The six-basis inverse mapping absorbs this factor of
$5$, restoring integer six-vector coordinates.  As a consequence, every
reflection is implementable as a precomputed $6\times 6$ integer matrix,
costing only integer multiplies and adds at runtime.

Finally we describe two GPU architectures built on this foundation:
a batch arithmetic engine (Mode~1) and a $\varphi$-native crystal
lattice simulator (Mode~2) in which each GPU thread is permanently
assigned to one lattice node and communicates only with its six neighbors.
Mode~2 constitutes a new computational model whose topology is the
quasicrystal itself.
\end{abstract}

\tableofcontents
\newpage

% ─────────────────────────────────────────────────────────────────────────────
\section{Motivation}

A slide rule works because logarithms convert multiplication into
addition.  The physical displacement of one scale against another
encodes the product.  The key insight is that the number line can be
re-indexed---by powers of $10$, or any convenient base---so that the
index does the arithmetic rather than the value itself.

The digital slide rule presented here works on the same principle,
with two differences.

\medskip
\noindent\textbf{First}, the base is $\varphi$ instead of $10$.
This choice is not arbitrary.  The golden ratio satisfies
\[
  \varphi^{2} = \varphi + 1,
\]
so every power of $\varphi$ is an integer linear combination of
$\varphi$ and $1$.  The number system is \emph{closed}: computation
never leaves the lattice.

\medskip
\noindent\textbf{Second}, the ``slide'' is a precomputed table of
integer pairs rather than a physical scale.  Lookup is an array-index
operation; advancing along the scale is integer addition on the index.
The entire computation runs in integer arithmetic from start to finish.

The system is especially well suited to the geometry of dodecahedral
and quasicrystal structures, where the natural angles---multiples of
$36^\circ$---produce sines and cosines that live \emph{exactly} in the
$\Zphi$ lattice.

% ─────────────────────────────────────────────────────────────────────────────
\section{The Number Representation}

\begin{definition}[PhiBase notation]
A \emph{PhiBase number} is an element of $\Zphi = \{p\varphi + n \mid
p,n\in\Z\}$.  We write it as $\Phib{p}{n}$ to denote the value
$p\varphi + n$.
\end{definition}

The \emph{positive cone} restricts to $p,n \geq 0$.  All lattice
node coordinates in Mode~2 (Section~\ref{sec:gpu}) lie in the
positive cone.

\medskip
Table~\ref{tab:decode} lists the seven coordinate symbols used throughout
the dodecahedral geometry code.  Every vertex of every Platonic solid,
and every vertex of the four Kepler--Poinsot solids, has coordinates
drawn from this vocabulary.

\begin{table}[h]
\centering
\caption{The seven PhiBase coordinate symbols and their values.}
\label{tab:decode}
\begin{tabular}{ccl}
\toprule
Symbol & $\Phib{p}{n}$ & Float value \\
\midrule
\texttt{z} & $\Phib{0}{0}$  & $0$ \\
\texttt{O} & $\Phib{0}{1}$  & $1$ \\
\texttt{o} & $\Phib{0}{-1}$ & $-1$ \\
\texttt{f} & $\Phib{-1}{1}$ & $1/\varphi \approx 0.618$ \\
\texttt{F} & $\Phib{1}{-1}$ & $-1/\varphi \approx -0.618$ \\
\texttt{p} & $\Phib{-1}{0}$ & $-\varphi \approx -1.618$ \\
\texttt{P} & $\Phib{1}{0}$  & $\varphi \approx 1.618$ \\
\bottomrule
\end{tabular}
\end{table}

\medskip
A few PhiBase values illustrate the Fibonacci structure:
\[
  \Phib{0}{1}=1,\quad
  \Phib{1}{0}=\varphi,\quad
  \Phib{1}{1}=\varphi^{2},\quad
  \Phib{2}{1}=\varphi^{3},\quad
  \Phib{3}{2}=\varphi^{4},\quad
  \Phib{5}{3}=\varphi^{5}.
\]
The coefficients $(p,n)$ at each step are consecutive Fibonacci numbers.
This is the central structural fact, proved in Section~\ref{sec:spine}.

For values less than $1$, larger coefficients with mixed signs reach
them in the full ring: for example $\Phib{5}{-8} = 5\varphi - 8
\approx 0.090$.

% ─────────────────────────────────────────────────────────────────────────────
\section{The Two Tables}

\subsection{Table~2 — The Fibonacci Spine}
\label{sec:spine}

\begin{definition}[Fibonacci Spine]
The \emph{Fibonacci Spine} is the sequence of entries
\[
  T_2[k] \;=\; \bigl(F_{k+1},\; F_k,\; \varphi^k\bigr),
  \qquad k = 0, 1, 2, \ldots
\]
where $F_k$ denotes the $k$-th Fibonacci number ($F_0=0,F_1=1$).
Entry $k$ represents the PhiBase number $\Phib{F_{k+1}}{F_k}$.
\end{definition}

\begin{theorem}[Fibonacci representation of $\varphi^k$]
\label{thm:fibpow}
For all $k \geq 0$,
\[
  \varphi^k \;=\; F_{k+1}\,\varphi + F_k.
\]
\end{theorem}

\begin{proof}
By induction.  Base cases: $\varphi^0 = 1 = 0\cdot\varphi + 1$
($F_1=1, F_0=0$ \checkmark) and $\varphi^1 = \varphi = 1\cdot\varphi+0$
($F_2=1,F_1=1$... adjusting: $F_2=1,F_1=1$, giving $1\cdot\varphi+0$\checkmark).
Inductive step: assume $\varphi^k = F_{k+1}\varphi + F_k$.  Then
\begin{align*}
  \varphi^{k+1} &= \varphi\cdot(F_{k+1}\varphi + F_k)
               = F_{k+1}\varphi^2 + F_k\varphi \\
               &= F_{k+1}(\varphi+1) + F_k\varphi
               = (F_{k+1}+F_k)\varphi + F_{k+1}
               = F_{k+2}\varphi + F_{k+1}. \qedhere
\end{align*}
\end{proof}

The spine has three immediate consequences for arithmetic.

\begin{corollary}[Integer arithmetic on the spine]
\label{cor:spineops}
Let $j,k \geq 0$.  Then:
\begin{align}
  \varphi^j \times \varphi^k &= \varphi^{j+k}
    &&\text{(index addition)} \label{eq:mul}\\
  \varphi^j / \varphi^k &= \varphi^{j-k}
    &&\text{(index subtraction, } j\geq k\text{)} \label{eq:div}\\
  \varphi^j > \varphi^k &\iff j > k
    &&\text{(index comparison)} \label{eq:cmp}
\end{align}
No floating-point operations are required.
\end{corollary}

\paragraph{The coupled step.}
Moving one step along a basis direction at spine level $k$ induces an
automatic displacement of $1/\varphi$ in the conjugate direction.
Since $\varphi^{-(1)} = F_2\varphi - F_1 = \varphi - 1 = 1/\varphi$,
this coupled displacement is exactly spine entry $k+1$.  Every spine
entry therefore carries a pointer to the next entry; no additional data
is stored.  This \emph{double-hit} property encodes the inflation rule
of the dodecahedral quasicrystal.

\medskip
Table~\ref{tab:spine} shows the first nine entries.  The spine contains
approximately 40 entries for $10^{-8}$ precision, since
$\varphi^{-38} \approx 10^{-8}$.

\begin{table}[h]
\centering
\caption{Fibonacci Spine entries $T_2[k]$, $k = 0,\ldots,8$.}
\label{tab:spine}
\begin{tabular}{rrrl}
\toprule
$k$ & $F_{k+1}$ & $F_k$ & $\varphi^k$ (float) \\
\midrule
0 & 1 & 0 & 1.000\,000 \\
1 & 1 & 1 & 2.618\,034 \\
2 & 2 & 1 & 4.236\,068 \\
3 & 3 & 2 & 6.854\,102 \\
4 & 5 & 3 & 11.090\,17 \\
5 & 8 & 5 & 17.944\,27 \\
6 & 13& 8 & 29.034\,44 \\
7 & 21& 13& 46.979\,00 \\
8 & 34& 21& 76.013\,16 \\
\bottomrule
\end{tabular}
\end{table}

\subsection{Table~1 — The Dense Grid}

The spine covers only the exact powers of $\varphi$.  The \emph{Dense
Grid} fills in all values between them.

\begin{definition}[Dense Grid]
Let $N$ be a positive integer bound.  The \emph{Dense Grid}
$T_1(N)$ is the sorted sequence of all distinct values
$p\varphi + n$ with $0 \leq p,n \leq N$, together with:
\begin{itemize}[nosep]
  \item their \emph{rank} (position in sorted order),
  \item their exact $(p,n)$ pair,
  \item their float value (precompute and display only), and
  \item their \emph{spine bracket} $k^*$, the largest $k$ with
        $\varphi^k \leq p\varphi+n$.
\end{itemize}
\end{definition}

For $N = 25$ the grid contains approximately $13{,}000$ entries and
achieves maximum approximation error below $10^{-8}$ across the unit
interval $[0,1]$.

The \emph{rank} of an entry is its working integer identity.  All
comparisons reduce to rank comparisons; all nearest-value queries use
binary search on the sorted float column, performed once at query time.
The float values are never used again after the precompute phase.

\subsection{The Slide Rule Analogy}

The classical slide rule has a fixed \emph{stock} and a sliding
\emph{rule}, both graduated logarithmically.  Multiplication is
implemented by physical displacement: adding logarithms.

The digital slide rule has the same two-layer structure:

\begin{center}
\renewcommand{\arraystretch}{1.3}
\begin{tabular}{ll}
\toprule
Slide rule & Digital slide rule \\
\midrule
Stock & Fibonacci Spine (Table~2) \\
Sliding rule & Dense Grid (Table~1) \\
Decade & Spine index $k$ \\
Mantissa position & Grid rank \\
Physical displacement & Index addition \\
\bottomrule
\end{tabular}
\end{center}

To multiply two PhiBase values: express each as a spine index plus
a grid rank, add the spine indices (Corollary~\ref{cor:spineops}),
combine the grid offsets, and read the result.  No actual multiplication
of real numbers occurs.

The precision is controlled by the grid size $N$: larger $N$ buys
finer coverage at the cost of table size.  At $N=25$ the system
achieves $10^{-8}$ precision with $\approx 13{,}000$ entries.

\subsection{Function Lookup}

Any function $f : D \to [0,1]$ can be encoded in the Dense Grid.
The function is evaluated once during precompute; at runtime only
the rank is used.

\begin{table}[h]
\centering
\caption{Sample function values encoded in the Dense Grid.}
\label{tab:funclookup}
\begin{tabular}{lrrl}
\toprule
Expression & True value & Grid value & Error \\
\midrule
$\sin(\pi/6)$   & 0.500\,000\,0 & 0.500\,000\,0 & $2.1\times10^{-9}$ \\
$\sin(\pi/4)$   & 0.707\,106\,8 & 0.707\,106\,8 & $1.4\times10^{-9}$ \\
$\sin(\pi/3)$   & 0.866\,025\,4 & 0.866\,025\,4 & $8.3\times10^{-10}$ \\
$\cos(\pi/5)$   & 0.809\,016\,9 & 0.809\,016\,9 & $0$ (exact) \\
$e^{-1}$        & 0.367\,879\,4 & 0.367\,879\,4 & $3.1\times10^{-9}$ \\
$1/\varphi$     & 0.618\,033\,9 & 0.618\,033\,9 & $0$ (exact) \\
\bottomrule
\end{tabular}
\end{table}

The zero errors for $\cos(\pi/5) = \varphi/2$ and $1/\varphi$ are
not coincidental.  The angles native to dodecahedral geometry produce
trigonometric values that are exact PhiBase numbers.  The system is
algebraically matched to the geometry it was designed for.

For values outside $[0,1]$, factor out the appropriate $\varphi^k$,
look up the mantissa in the unit grid, and record the pair $(k,
\text{rank})$.  Two integers represent any positive real to $10^{-8}$
precision.

% ─────────────────────────────────────────────────────────────────────────────
\section{Why This Is a Win}

Classical IEEE~754 floating-point arithmetic accumulates rounding error
at every operation.  In long computation chains---building a crystal
lattice, propagating a wave, computing structural stress---accumulated
error is a practical problem.  The PhiBase system eliminates it.

\begin{enumerate}[leftmargin=*]
\item \textbf{No runtime rounding error.}  Every operation during
computation is on integers.  The only approximation occurs at the
initial encoding step, where a real number is matched to its nearest
grid entry.  That error is bounded by $10^{-8}$ and does not grow.
It is committed once and frozen.

\item \textbf{Exact membership testing.}  Asking whether a computed
point belongs to the lattice is ill-posed in floating point.  In
PhiBase it is exact: either the result is a valid $(p,n)$ pair or
it is not.  This enables exact verification of geometric membership,
a prerequisite for the GPU crystal lattice (Section~\ref{sec:gpu}).

\item \textbf{Universal function encoding.}  Any function with range
in $[0,1]$ maps into the Dense Grid.  This includes all trigonometric
functions restricted to standard quarter-periods, all probability
distributions, and all functions arising from the icosahedral geometry.
The table is built once; thereafter, function evaluation is an integer
array lookup.
\end{enumerate}

The system is not universally faster than floating point.  A single
lookup carries array-indexing overhead.  The win is in correctness and
in the ability to run the entire geometric engine on integer hardware,
including GPU thread indices, without a floating-point unit.

% ─────────────────────────────────────────────────────────────────────────────
\section{The Six-Basis Lattice and Icosahedral Geometry}

\subsection{The Six Basis Vectors}

The dodecahedron and icosahedron share the same symmetry group
$I_h$ of order $120$.  Their vertices and edge midpoints span six
distinct directions, which we take as basis vectors:

\[
  \mathbf{B}_1 = (\varphi,1,0),\quad
  \mathbf{B}_2 = (\varphi,-1,0),\quad
  \mathbf{B}_3 = (1,0,\varphi),
\]
\[
  \mathbf{B}_4 = (-1,0,\varphi),\quad
  \mathbf{B}_5 = (0,\varphi,1),\quad
  \mathbf{B}_6 = (0,\varphi,-1).
\]

Every lattice point is a six-vector $[a,b,c,d,e,f]\in\Z^6$, mapping
to Cartesian $\Zphi^3$ coordinates via:
\begin{align}
  x &= \varphi(a-b) + (e-f), \label{eq:x}\\
  y &= (c-d) + \varphi(e+f), \label{eq:y}\\
  z &= (a+b) + \varphi(c+d). \label{eq:z}
\end{align}

In Mode~2 (Section~\ref{sec:gpu}), all six components are non-negative.
Neighbors are reached by incrementing one component by one---a purely
additive operation.

\subsection{The Fifteen Symmetry Planes}

The icosahedral reflection group $I_h$ has exactly $15$ symmetry planes,
arranged in five sets of three mutually orthogonal planes.  Each plane
contains two opposite edges of the icosahedron and two opposite edges
of the dodecahedron.

Six of the fifteen planes are the face-normal planes of the dodecahedron,
which we label $A$ through $F$.  Their normal vectors in $\Zphi^3$ are:

\begin{align*}
  \mathbf{n}_A &= (\varphi,\; 1,\; 0) = (\Phib{1}{0},\Phib{0}{1},\Phib{0}{0}), \\
  \mathbf{n}_B &= (\varphi,\;-1,\; 0) = (\Phib{1}{0},\Phib{0}{-1},\Phib{0}{0}), \\
  \mathbf{n}_C &= (1,\; 0,\; \varphi) = (\Phib{0}{1},\Phib{0}{0},\Phib{1}{0}), \\
  \mathbf{n}_D &= (-1,\; 0,\;\varphi) = (\Phib{0}{-1},\Phib{0}{0},\Phib{1}{0}), \\
  \mathbf{n}_E &= (0,\;\varphi,\; 1) = (\Phib{0}{0},\Phib{1}{0},\Phib{0}{1}), \\
  \mathbf{n}_F &= (0,\;\varphi,\;-1) = (\Phib{0}{0},\Phib{1}{0},\Phib{0}{-1}).
\end{align*}

The remaining nine planes are generated by composing these six.  In
the implementation, a point reflected through planes $A$ then $B$ is
denoted with the suffix \texttt{\textasciitilde AB}, building up
words in the reflection group.  The six generators $A$--$F$ generate
the full group $I_h$.

% ─────────────────────────────────────────────────────────────────────────────
\section{The Integrality Theorem for Reflections}
\label{sec:reflect}

\subsection{The Reflection Formula}

Reflection of a point $\mathbf{p}\in\Zphi^3$ through the plane with
normal $\mathbf{n}$ is:
\[
  \mathbf{p}' = \mathbf{p} - 2\,\frac{\mathbf{p}\cdot\mathbf{n}}
                                      {\mathbf{n}\cdot\mathbf{n}}\,\mathbf{n}.
\]
The key quantity is $\mathbf{n}\cdot\mathbf{n}$.

\subsection{The Denominator is Always 5}

\begin{lemma}[Universal plane norm]
\label{lem:norm}
For all six planes $A$ through $F$,
\[
  \mathbf{n}\cdot\mathbf{n} \;=\; \Phib{1}{2} \;=\; \varphi + 2.
\]
\end{lemma}

\begin{proof}
By direct computation for plane $A$:
\[
  \mathbf{n}_A\cdot\mathbf{n}_A
  = \varphi^2 + 1^2 + 0^2
  = (\varphi+1) + 1
  = \varphi + 2
  = \Phib{1}{2}.
\]
All six normals share the same structure: two non-zero coordinates of
magnitude $\varphi$ or $1$, one zero coordinate.  Hence
$\mathbf{n}\cdot\mathbf{n} = \varphi^2 + 1 = \varphi + 2$ in every case.
\end{proof}

\begin{lemma}[Algebraic norm]
\label{lem:algnorm}
The algebraic norm of $\Phib{1}{2} = \varphi+2$ is $5$.
\end{lemma}

\begin{proof}
For $\alpha = p\varphi + n \in \Zphi$, the algebraic norm is
$\mathcal{N}(\alpha) = n^2 + np - p^2$.
For $\Phib{1}{2}$: $\mathcal{N} = 4 + 2 - 1 = 5$.
\end{proof}

\begin{corollary}
\label{cor:denom5}
Division by $\mathbf{n}\cdot\mathbf{n}$ in $\Zphi$ is equivalent to
multiplying by $\Phib{-1}{3}/5$, since
\[
  \Phib{1}{2}\times\Phib{-1}{3} = \Phib{0}{5} = 5.
\]
Therefore the scale factor $\mathbf{p}\cdot\mathbf{n}/(\mathbf{n}\cdot\mathbf{n})$
lies in $\tfrac{1}{5}\Zphi$, and every reflected coordinate lies in
$\tfrac{1}{5}\Zphi$.
\end{corollary}

\subsection{The Main Theorem}

\begin{theorem}[Reflections preserve the lattice]
\label{thm:reflect}
Let $\mathbf{v}\in\Z^6$ be any six-basis lattice vector.  Let $R_X$
denote reflection through any of the six face-normal planes $X\in
\{A,B,C,D,E,F\}$.  Then $R_X(\mathbf{v})$ is also a six-basis lattice
vector with integer components.
\end{theorem}

\begin{proof}[Proof sketch]
\begin{enumerate}[label=\arabic*.,leftmargin=*]
\item \textbf{Forward map.}  Apply equations~\eqref{eq:x}--\eqref{eq:z}
to obtain Cartesian coordinates $(x,y,z)\in\Zphi^3$.

\item \textbf{Reflect.}  By Lemma~\ref{lem:norm} and
Corollary~\ref{cor:denom5}, the reflected coordinates
$(x',y',z')$ lie in $\tfrac{1}{5}\Zphi^3$, with denominator
exactly $5$.

\item \textbf{Inverse map.}  Apply the inverse of
equations~\eqref{eq:x}--\eqref{eq:z}.  The six-basis inverse map
is a $\Z$-linear combination of the Cartesian coordinates.  The
factor of $5$ in the denominator cancels exactly against the
structure of the dodecahedral map, yielding integer components
$a',b',c',d',e',f'\in\Z$.
\end{enumerate}
\end{proof}

\begin{corollary}[Group action on the lattice]
The six reflections $R_A,\ldots,R_F$ generate a group of symmetries
acting on the six-basis lattice.  Every orbit consists entirely of
valid lattice points.  The lattice is closed under the full icosahedral
reflection group $I_h$.
\end{corollary}

\subsection{The Twelve Exact Cases}

When $\mathbf{p}\cdot\mathbf{n} = 0$, the point lies on the reflection
plane and the reflection is the identity.  Numerical verification shows
that each of the twelve icosahedron vertices lies on exactly one of
the six planes $A$--$F$.  These are the only twelve cases in which
the reflected coordinates lie in $\Zphi$ directly (rather than
$\tfrac{1}{5}\Zphi$).

\subsection{Sample Reflected Coordinates}

Table~\ref{tab:reflect} shows computed reflected values for a selection
of icosahedron vertices.  Coordinates in $\tfrac{1}{5}\Zphi$ are
written $(a\varphi+b)/5$.

\begin{table}[h]
\centering
\caption{Sample reflections of icosahedron vertices.}
\label{tab:reflect}
\begin{tabular}{lllll}
\toprule
Point & Plane & $x'$ & $y'$ & $z'$ \\
\midrule
$(0,1,-1/\varphi)$ & $A$ & $\frac{-4\varphi+2}{5}$ & $\frac{2\varphi-1}{5}$ & $\frac{5\varphi-5}{5}$ \\[4pt]
$(0,1,-1/\varphi)$ & $B$ & $\frac{4\varphi-2}{5}$  & $\frac{2\varphi-1}{5}$ & $\frac{5\varphi-5}{5}$ \\[4pt]
$(0,1,-1/\varphi)$ & $C$ & $\frac{2\varphi-6}{5}$  & $\frac{5}{5}$          & $\frac{\varphi-3}{5}$  \\[4pt]
$(0,-1,-1/\varphi)$& $A$ & $\frac{4\varphi-2}{5}$  & $\frac{-2\varphi+1}{5}$& $\frac{5\varphi-5}{5}$ \\
\bottomrule
\end{tabular}
\end{table}

Note: $(5\varphi-5)/5 = \varphi - 1 = 1/\varphi = \Phib{-1}{1}$, so
some coordinates simplify back into $\Zphi$.

\subsection{GPU Implementation of Reflections}

The three-step procedure---forward map, reflect, inverse map---can be
precomputed as a single $6\times 6$ integer matrix $M_X$ for each
plane $X$.  The matrix entries are integers because the factor of $5$
is absorbed symbolically during precomputation.

At runtime, reflecting a lattice node costs:
\[
  \mathbf{v}' = M_X\,\mathbf{v},
  \qquad M_X \in \Z^{6\times 6},
\]
which requires exactly 36 integer multiplications and 30 integer
additions---or, exploiting sparsity, substantially fewer.  No
floating-point unit is needed.

\begin{table}[h]
\centering
\caption{Summary of reflection integrality properties.}
\label{tab:reflsummary}
\begin{tabular}{ll}
\toprule
Property & Value \\
\midrule
Plane norm $\mathbf{n}\cdot\mathbf{n}$ & $\Phib{1}{2} = \varphi+2$ \\
Algebraic norm of $\Phib{1}{2}$ & $5$ \\
Reflected coords live in & $\tfrac{1}{5}\Zphi$ \\
Denominator (all planes) & always $5$ \\
Exact cases (point on plane) & $12$ (one per icosahedron vertex) \\
Six-basis absorbs the $/5$ & yes — lattice preserved \\
GPU implementation & $6\times 6$ integer matrix \\
\bottomrule
\end{tabular}
\end{table}

% ─────────────────────────────────────────────────────────────────────────────

\section{The Inverse Map and Lattice Operations}
\label{sec:inverse}

\subsection{The Inverse Map Explicitly}

The forward map from six-vector $[a,b,c,d,e,f]\in\Z^6$ to Cartesian
$\mathbb{Z}[\varphi]^3$ is given by equations~\eqref{eq:x}--\eqref{eq:z}.
Separating the $\varphi$-coefficients from the integer parts:

\begin{align}
  x = P(p_x, n_x) &\quad\text{where}\quad p_x = a-b,\quad n_x = e-f \label{eq:xpn}\\
  y = P(p_y, n_y) &\quad\text{where}\quad p_y = e+f,\quad n_y = c-d \label{eq:ypn}\\
  z = P(p_z, n_z) &\quad\text{where}\quad p_z = c+d,\quad n_z = a+b \label{eq:zpn}
\end{align}

This gives six linear equations in six unknowns.  Solving:

\begin{align}
  a &= \tfrac{1}{2}(n_z + p_x), &
  b &= \tfrac{1}{2}(n_z - p_x), \label{eq:ab}\\
  c &= \tfrac{1}{2}(p_z + n_y), &
  d &= \tfrac{1}{2}(p_z - n_y), \label{eq:cd}\\
  e &= \tfrac{1}{2}(p_y + n_x), &
  f &= \tfrac{1}{2}(p_y - n_x). \label{eq:ef}
\end{align}

\subsection{The Parity Constraint}

The factor of $\tfrac{1}{2}$ in equations~\eqref{eq:ab}--\eqref{eq:ef}
imposes a necessary condition for a Cartesian PhiBase triple to be the
image of an integer six-vector.

\begin{definition}[Lattice parity]
A Cartesian triple $P(p_x,n_x),\,P(p_y,n_y),\,P(p_z,n_z) \in
\mathbb{Z}[\varphi]^3$ is \emph{lattice-compatible} if all six
quantities
\[
  n_z \pm p_x,\quad p_z \pm n_y,\quad p_y \pm n_x
\]
are even integers.
\end{definition}

\begin{theorem}[Inverse map integrality]
\label{thm:inverse}
A Cartesian triple $(x,y,z) \in \mathbb{Z}[\varphi]^3$ is the image
of an integer six-vector under the forward map if and only if it is
lattice-compatible.  When it is, the six-vector is uniquely determined
by equations~\eqref{eq:ab}--\eqref{eq:ef}.
\end{theorem}

Numerical verification confirms that all six-vectors arising from
integer inputs satisfy the forward map and recover exactly under the
inverse map.

\subsection{The Full Round Trip for Reflections}

Combining the reflection result (Section~\ref{sec:reflect}) with the
inverse map, the full round trip is:

\[
  [a,b,c,d,e,f]
  \;\xrightarrow{\text{forward}}\;
  (x,y,z) \in \mathbb{Z}[\varphi]^3
  \;\xrightarrow{\text{reflect}}\;
  (x',y',z') \in \tfrac{1}{5}\mathbb{Z}[\varphi]^3
  \;\xrightarrow{\text{inverse}}\;
  [a',b',c',d',e',f']
\]

Writing the reflected coordinates as $(a_x\varphi + b_x)/5$ etc., the
inverse map gives:
\[
  a' = \frac{b_z + a_x}{10}, \quad
  b' = \frac{b_z - a_x}{10}, \quad
  c' = \frac{a_z + b_y}{10}, \quad
  d' = \frac{a_z - b_y}{10}, \quad
  e' = \frac{a_y + b_x}{10}, \quad
  f' = \frac{a_y - b_x}{10}.
\]

For the round trip to yield integers, each numerator must be divisible
by $10 = 2 \times 5$.  The factor of $5$ comes from the reflection
denominator (Lemma~\ref{lem:algnorm}); the factor of $2$ comes from
the inverse map parity constraint (Theorem~\ref{thm:inverse}).

\begin{theorem}[Reflection round-trip integrality]
\label{thm:roundtrip}
For any six-vector $[a,b,c,d,e,f]$ whose Cartesian coordinates are
lattice-compatible, the reflected six-vector $[a',b',c',d',e',f']$
has integer components.  The denominator in the full round trip is
exactly $10$.
\end{theorem}

This strengthens Theorem~\ref{thm:reflect}: the lattice is closed
under reflection, and the precise arithmetic cost of reflection is
a single division by $10$ at the end of the precomputed $6\times 6$
matrix computation.

\subsection{Translation in Six-Basis}

Translation by one step in the $i$-th basis direction increments one
component of the six-vector by one:
\[
  [a,b,c,d,e,f] \;\longmapsto\; [a,b,c,d,e,f] + \mathbf{e}_i.
\]

In Cartesian terms this adds the corresponding basis vector, which is
already in $\mathbb{Z}[\varphi]^3$.  No division occurs; the parity
constraint is preserved.  Translation is the cheapest possible lattice
operation: one integer increment.

This is why the Mode~2 GPU lattice (Section~\ref{sec:gpu}) uses
translation as its fundamental step.  Each clock tick increments
neighbor relationships by one six-vector component---nothing more.

\subsection{Phi-Expansion in Six-Basis}

\emph{Phi-expansion} scales all Cartesian coordinates by $\varphi$,
enlarging the structure by the golden ratio.  Since
$\varphi \cdot P(p,n) = P(p+n,\,p)$ (from $\varphi^2=\varphi+1$),
the expanded Cartesian coordinates are:

\begin{align*}
  \varphi\,x &= P(p_x + n_x,\; p_x), \\
  \varphi\,y &= P(p_y + n_y,\; p_y), \\
  \varphi\,z &= P(p_z + n_z,\; p_z).
\end{align*}

Substituting equations~\eqref{eq:xpn}--\eqref{eq:zpn} and applying
the inverse map:

\begin{align}
  a' &= \tfrac{1}{2}\bigl((c+d) + (a-b)+(e-f)\bigr), \label{eq:phia}\\
  b' &= \tfrac{1}{2}\bigl((c+d) - (a-b)-(e-f)\bigr), \\
  c' &= \tfrac{1}{2}\bigl((c+d)+(a+b) + (e+f)\bigr), \\
  d' &= \tfrac{1}{2}\bigl((c+d)+(a+b) - (e+f)\bigr), \\
  e' &= \tfrac{1}{2}\bigl((e+f)+(c-d) + (a-b)-(e-f)\bigr), \\
  f' &= \tfrac{1}{2}\bigl((e+f)+(c-d) - (a-b)+(e-f)\bigr). \label{eq:phif}
\end{align}

\begin{theorem}[Phi-expansion parity condition]
\label{thm:phiexp}
Phi-expansion of a six-vector $[a,b,c,d,e,f]$ yields an integer
six-vector if and only if all six quantities in
\eqref{eq:phia}--\eqref{eq:phif} are even.

A sufficient condition is that the components satisfy the
\emph{even-sum parity}:
\[
  (a+b+c+d+e+f) \equiv 0 \pmod{2}.
\]
Not all six-vectors satisfy this; phi-expansion is valid on the
sublattice of even-sum vectors.
\end{theorem}

Numerical verification shows that individual basis vectors
$\mathbf{e}_i$ do \emph{not} satisfy the parity condition, but pairs
such as $[1,1,0,0,0,0]$ and $[1,0,1,0,0,0]$ do.  This confirms
that phi-expansion acts on pairs of basis steps, consistent with
the double-hit coupling property of the spine (Section~\ref{sec:spine}).

\begin{remark}
The three operations --- translation, phi-expansion, and reflection
--- have a clear cost hierarchy in the six-basis:

\begin{center}
\renewcommand{\arraystretch}{1.3}
\begin{tabular}{llll}
\toprule
Operation & Division & Parity condition & GPU cost \\
\midrule
Translation   & none  & always valid      & 1 integer increment \\
Phi-expansion & by 2  & even-sum sublattice & 6 adds, 1 right-shift \\
Reflection    & by 10 & lattice-compatible & $6\times6$ integer matrix \\
\bottomrule
\end{tabular}
\end{center}

Translation is free.  Phi-expansion costs one bit-shift.  Reflection
costs a matrix multiply with a compile-time divisor of $10$.
\end{remark}

\subsection{Chaining Operations}

Since all three operations return integer six-vectors (subject to the
stated parity conditions), they can be chained freely:

\[
  \text{translate} \to \text{reflect} \to \text{phi-expand}
  \to \text{translate} \to \cdots
\]

Each step takes integer input and produces integer output.  The
computation never leaves the lattice.  This is the key property
that makes the six-basis representation suitable for the GPU crystal
lattice: every operation in the geometric engine runs on integers from
start to finish, with no floating-point arithmetic required at any step.

\section{The GPU Crystal Lattice}
\label{sec:gpu}

\subsection{Mode 1 — GPU as Arithmetic Engine}

The first mode uses the GPU as a batch arithmetic engine serving the
CPU.  The CPU submits geometric problems---multiply two PhiBase numbers,
find the nearest lattice point to a real value, evaluate a trigonometric
function---and the GPU solves them in parallel, one per thread.

The Dense Grid and Fibonacci Spine reside in GPU constant memory.
A batch of $N$ queries dispatches $N$ threads simultaneously; each
thread performs one array lookup and returns an integer rank.  No
floating-point arithmetic occurs in the kernel.

The throughput benefit depends on batch size.  A single query is slower
on the GPU due to dispatch overhead.  Ten thousand queries in parallel
is where the GPU wins.  In geometric applications---all vertices of a
structure, all edge lengths, all angles---batches of this size arise
naturally.

Mode~1 does not change what is being computed; it changes where and
how fast.

\subsection{Mode 2 — The $\varphi$-Native Crystal Computer}

Mode~2 is different in kind, not just degree.

Each GPU thread is permanently assigned to one lattice node.  The
node's address is a six-vector $[a,b,c,d,e,f]$ with all non-negative
integer components (the positive cone).  The thread knows its six
neighbors, reached by incrementing one component by one.  This
adjacency list is fixed at initialization and never changes.

Each thread maintains two state values: \texttt{last} and \texttt{next}.
At each clock tick the thread:
\begin{enumerate}[nosep]
  \item reads \texttt{last} from all six neighbors,
  \item computes its own \texttt{next} from those values, and
  \item waits at a synchronization barrier.
\end{enumerate}
When all threads have finished, the barrier releases, \texttt{next}
becomes \texttt{last}, and the cycle repeats.

The CPU's only roles are loading initial conditions, starting the clock,
and reading final state.  Between clock ticks the CPU is idle.
The lattice runs itself.

The Metal shader kernel is:

\begin{lstlisting}[language=C]
kernel void lattice_step(
    device PhiNode* last  [[ buffer(0) ]],
    device PhiNode* next  [[ buffer(1) ]],
    device int*     adj   [[ buffer(2) ]],
    uint   id             [[ thread_position_in_grid ]])
{
    int n0=adj[id*6+0], n1=adj[id*6+1], n2=adj[id*6+2],
        n3=adj[id*6+3], n4=adj[id*6+4], n5=adj[id*6+5];

    next[id].value = update(
        last[n0], last[n1], last[n2],
        last[n3], last[n4], last[n5]);
}
\end{lstlisting}

The adjacency buffer \texttt{adj} is precomputed from the six-basis
arithmetic and uploaded once.  The kernel itself contains only integer
array reads, the \texttt{update} function, and a write.

\subsection{Why Mode 2 Is a New Computational Model}

A CPU cannot perform this computation, not because of speed but because
the computation \emph{is} the simultaneous local interaction of every
node with its six neighbors.  That has no sequential equivalent.

The update function is left open deliberately.  It could encode:
\begin{itemize}[nosep]
  \item a wave equation (acoustic or electromagnetic propagation),
  \item a stress or strain field computation,
  \item a signal routing or switching rule, or
  \item a physics simulation native to quasicrystal geometry.
\end{itemize}

The lattice provides the geometry.  The update function provides the
physics.  Together they define a programmable quasicrystal computer
whose wiring diagram is the dodecahedral lattice itself.

A signal propagating through this lattice travels along $\varphi$-scaled
paths and reflects off $\varphi$-symmetric boundaries.  The propagation
properties are those of physical quasicrystals---directional,
non-dispersive along the icosahedral axes---arising here not from
material science but from the algebraic structure of the lattice.

\subsection{Comparison of the Two Modes}

\begin{center}
\renewcommand{\arraystretch}{1.3}
\begin{tabular}{lll}
\toprule
Property & Mode 1 & Mode 2 \\
\midrule
Role of GPU & serves CPU & \emph{is} the computation \\
Table used & Full $\Zphi$ & Positive cone only \\
Arithmetic & Signed & Non-negative \\
CPU during run & collects results & idle \\
Speedup question & yes (batch size) & not applicable \\
Novelty & engineering & new model \\
\bottomrule
\end{tabular}
\end{center}

Mode~1 is useful.  Mode~2 is novel.

% ─────────────────────────────────────────────────────────────────────────────
\section{Summary and Future Work}

We have presented four interlocking results.

\medskip
\noindent\textbf{1. The digital slide rule.}
PhiBase notation $\Phib{p}{n} = p\varphi + n$ with integer $p,n$
supports exact arithmetic via two precomputed tables.  The Fibonacci
Spine reduces multiplication to index addition; the Dense Grid reduces
comparison to rank comparison.  Any function with range in $[0,1]$
is encoded at $10^{-8}$ precision.  No floating-point arithmetic is
required during computation.

\medskip
\noindent\textbf{2. The integrality theorem.}
Reflection through any of the six dodecahedral face-normal planes maps
$\Zphi^3$ into $\tfrac{1}{5}\Zphi^3$, with denominator exactly $5$
in all cases (Theorem~\ref{thm:reflect}).  The six-basis inverse mapping
absorbs this factor, restoring integer lattice coordinates.  Reflections
therefore reduce to $6\times 6$ integer matrix multiplication at runtime.

\medskip
\noindent\textbf{3. The $\varphi$-native GPU lattice.}
The dodecahedral six-basis maps directly onto GPU thread space.
Mode~1 provides batch arithmetic acceleration.  Mode~2 creates a
new computational model---a programmable quasicrystal computer---in
which the lattice topology is the wiring diagram and each GPU thread
is a permanently assigned lattice node.

\medskip
Future work includes:

\begin{enumerate}[leftmargin=*]
\item Explicit derivation of the six $6\times 6$ reflection matrices
      and verification that they generate $I_h$.
\item Implementation of the Metal Mode~2 kernel and measurement of
      propagation properties for wave and stress update functions.
\item Extension to the full $\Zphi$ ring (signed coefficients) for
      Mode~1 table generation, with analysis of the required table
      depth for target precision.
\item Investigation of whether the $\tfrac{1}{5}\Zphi$ intermediate
      in reflection has physical significance in quasicrystal physics.
\end{enumerate}

% ─────────────────────────────────────────────────────────────────────────────
\appendix
\section{Arithmetic in $\mathbb{Z}[\varphi]$}

Multiplication in $\Zphi$ uses $\varphi^2 = \varphi+1$:
\begin{align*}
  (p_1\varphi+n_1)(p_2\varphi+n_2)
  &= p_1 p_2\varphi^2 + (p_1 n_2 + n_1 p_2)\varphi + n_1 n_2 \\
  &= p_1 p_2(\varphi+1) + (p_1 n_2 + n_1 p_2)\varphi + n_1 n_2 \\
  &= (p_1 p_2 + p_1 n_2 + n_1 p_2)\,\varphi + (p_1 p_2 + n_1 n_2).
\end{align*}

Division uses the conjugate $\overline{p\varphi+n} = -p\varphi+(p+n)$:
\[
  \frac{p_1\varphi+n_1}{p_2\varphi+n_2}
  = \frac{(p_1\varphi+n_1)\overline{(p_2\varphi+n_2)}}
         {(p_2\varphi+n_2)\overline{(p_2\varphi+n_2)}}
  = \frac{(p_1\varphi+n_1)(-p_2\varphi+p_2+n_2)}
         {\mathcal{N}(p_2\varphi+n_2)},
\]
where $\mathcal{N}(p\varphi+n) = n^2+np-p^2 \in \Z$.
For $\mathbf{n}\cdot\mathbf{n} = \Phib{1}{2}$, we have
$\mathcal{N} = 5$, establishing Lemma~\ref{lem:algnorm}.

\section{The Decode Vocabulary in Context}

The seven symbols of Table~\ref{tab:decode} span coordinates
$\{0, \pm1, \pm\varphi, \pm1/\varphi\}$.  These are the only values
that appear as vertex coordinates in any of the nine regular polyhedra
(five Platonic solids and four Kepler--Poinsot solids) built from golden
triangles.  The vocabulary is complete for this geometry.

\end{document}
